\documentclass[11pt]{article}
\usepackage[utf8]{inputenc}
\usepackage{amsmath,amssymb,amsthm}
\usepackage[margin=1in]{geometry}
\usepackage{hyperref}
\usepackage{xcolor}

\newcommand{\tideref}[2][]{\href{https://therisensea.org/tides/#2/}{\texttt{#2}}}
\emergencystretch=1.5em

\title{Gaussian covariance rigidity\\
(germbij tensor programme, opening tide J2)}
\author{Tide loop (Claude Fable 5, with GPT-5.6 Sol deliberation)}
\date{2026-08-09}

\begin{document}
\maketitle

\begin{abstract}
\noindent The degree-$k$ tensor-recovery programme opens with its
injectivity engine: a continuous homogeneous function of positive
degree and polynomial growth whose self-covariance against the
quadratic Gaussian vanishes is identically
zero.\footnote{\texttt{homogeneous\_eq\_zero\_of\_gaussianCovariance\_self\_eq\_zero}
in \texttt{Laplace/Multi/GaussianCovariance.lean}.}
This is the no-Isserlis idea from the original multivariate scoping
consult, now in Lean: no Wick pairings, no moment matrices, no
monomial enumeration --- full support of the Gaussian weight (a
direct ball argument) plus the variance identity do all the work.
The fresh scoping consult staged the whole programme J0--J7 and
ruled this tide the minimal good opener; its J0/J1 slices
(polynomial-growth integrability, continuous full-support rigidity)
are proven inline and are reusable at every recovery rung.
\end{abstract}

\section{Setting}

With the H-recovery milestone complete (PRs \#66--\#73), the germbij
note's next multivariate commission is recovering the higher Taylor
data $L_k$, $k \ge 3$, from higher normalized moments. The scoping
consult staged this J0--J7: growth-integrability adapter, Gaussian
full support, covariance rigidity, multilinear polarization, radial
Taylor coefficients, the rate-sensitive pairwise DCT (flagged
largest), single-degree recovery, and the induction. It ruled J2 the
right first tide: central, self-contained, independent of the hard
DCT, with an unambiguous acceptance criterion.

\section{The thing formalised}

With $\mathbb{E}_H[f] = \int f K_H / \int K_H$ and
$\mathrm{Cov}_H(f,g) = \mathbb{E}_H[fg] - \mathbb{E}_H[f]\,
\mathbb{E}_H[g]$: for $Q$ continuous, of polynomial growth, and
homogeneous of degree $k > 0$,
\[
\mathrm{Cov}_H(Q, Q) = 0 \implies Q = 0,
\]
together with the supporting API: `HasPolynomialGrowth` with product
and shift closure, integrability of every polynomial-growth
observable against $K_H$, and the full-support rigidity lemma
(continuous nonnegative, zero weighted integral $\Rightarrow$
identically zero).

\section{How the proof works}

The variance identity
$\int (Q - \mathbb{E}Q)^2 K_H = Z_H \cdot \mathrm{Cov}_H(Q,Q)$
is an expand-and-divide computation whose integrability side
conditions come from the growth closure. Vanishing self-covariance
therefore kills the weighted integral of the squared deviation; the
deviation is continuous and nonnegative, so full support forces
$Q \equiv \mathbb{E}Q$. Full support itself is the consult's direct
ball argument: at a point where the continuous integrand $f K_H$ is
positive, continuity gives a ball where it exceeds half its value,
and the ball has positive measure --- no normalized measure, no
`Measure.support` API. Finally a constant homogeneous function of
positive degree vanishes, since $Q(0) = 0^k Q(0) = 0$.

\section{Where it stalled and how it moved}

No stalls; three quick passes, all catalogued classes (the
`Pi.sub` single-lambda witness for \texttt{integral\_add}, a
\texttt{set}-folded hypothesis making a rewrite dead) plus one new
niggle for the catalogue: this Mathlib pin deprecates
\texttt{push\_neg} in favour of \texttt{push Not}. The consult's
pre-flight warning --- check that the existing kernel-moment theorem
really is the norm-moment form, not powers of the kernel --- was
verified (it is) and saved a potential wrong turn.

\section{Mathlib lookup versus new mathematics}

Lookup: ball positivity, the set-integral comparisons, and the H2
integrability theorems.\footnote{\texttt{measure\_ball\_pos},
\texttt{setIntegral\_ge\_of\_const\_le},
\texttt{setIntegral\_le\_integral}.} New (project-level): the growth certificate and its
closure algebra, the weighted-integral expectation/covariance
definitions (deliberately not the probability-theory variance API),
the full-support rigidity lemma, and the homogeneous injectivity
theorem that every later recovery rung will instantiate at
$P = Q$.

\section{What this opens}

J3 (multilinear polarization: a symmetric $k$-linear form is
determined by its diagonal --- the consult expects a local
polarization lemma over subset sums), then J4 (radial Taylor
coefficients at degree $k$, cubic first), J5 (the rate-sensitive
pairwise DCT, the programme's flagged largest stage), and the
recovery ladder J6--J7.

\end{document}
